\documentclass[11pt,a4paper]{article}
\usepackage[margin=24mm]{geometry}
\usepackage{amsmath,amssymb,booktabs,longtable,array,hyperref}
\hypersetup{hidelinks}
\newcommand{\Lc}{\textbf{(c)}\ }
\newcommand{\Ld}{\textbf{(d)}\ }
\newcommand{\scr}{\texttt{audit\_\allowbreak phase2i.py}}
\newcolumntype{L}[1]{>{\raggedright\arraybackslash}p{#1}}
\emergencystretch=3em
\title{Independent audit of Phase 2i:\\normalisations, numerical evidence and unresolved identifications}
\author{Quantum Boussinesq solitons project}
\date{4 November 2026}
\begin{document}
\maketitle
\begin{abstract}
\Lc\ Phase 2i contains useful new conditional asymptotic calculations and correct
orientation bookkeeping, but its printed Wiener--Hopf theorem is incorrectly
normalised, its exchange criterion is false, and its assertion that all six
audits are resolved is unsupported. The original checker passes despite these
problems. Independent inverse Laplace transforms support the stated momentum
coefficients after a missing normalisation is restored. Independent phonon
quadratures confirm the orientation formulas, not their interpretation as a
physical determinant normalisation. \Ld\ The physical quantum frame and the
boundary-compatible projected determinant remain unestablished.
\end{abstract}

\section{Scope and reproducibility}
Labels retain their repository meanings: (a) established in literature, (b)
literature heuristic, (c) project derivation or computational evidence, (d)
unestablished interpretation or research judgment. All new mathematical checks
below are in \texttt{scripts/}\scr; numerical evidence is explicitly distinguished
from a proof. No existing file was edited, no git operation performed, and no
regularisation, physical charge or phase convention selected. The good-equation
conventions and the existing Toda extraction are retained. The independent
quantum-frame investigation is outside this audit's scope.

\Lc\ The Phase 2i source and PDF, the six named audits, the Phase 2f/2h definitions,
and relevant checkers were read. The Phase 2h independent audit was read in full
before implementation. This checker imports no Phase 2i formulas or audit
checker. Only the expressly permitted \texttt{FastNodesEdgeBA} finite-$h$ solver
is reused; its endpoint is set to exactly $q$, and its root bracket is expanded
when necessary. Analytical and integral checks use 60-digit mpmath; convolution
quadrature uses 35 digits; the dense Bethe solver and data regression use
float64. JSON stores raw data, input hashes, fit variants and timings. Hashes
are checked again at the end. Finite-$h$ results are cached; use
\texttt{python3 audit\_phase2i.py --refresh} to regenerate them.

The terms \emph{confirmed}, \emph{qualified}, and \emph{refuted} refer to the
specific claim, rather than the entire programme. In particular, ``confirmed
conditional algebra'' does not certify a physical S-matrix. \Ld

\section{Theorem 1: three distinct normalisation problems}
\paragraph{Fourier symbol: refuted as printed (\scr, \texttt{wh}).}
\Lc\ Write $T(s)=\widetilde\Phi_{\rm WH}(s)$ with Phase 2i's factor. For
$\omega\ne0$, direct multiplication gives
\begin{equation}
 \Phi_{\rm WH}(i\omega)\Phi_{\rm WH}(-i\omega)
 =\frac{\pi e^{|\omega|/2}}{\sinh(|\omega|/2)}
 =\frac{2\pi}{1-e^{-|\omega|}},
 \label{eq:symbol}
\end{equation}
not $2\pi^2 e^{\pi|\omega|}/\sinh(\pi|\omega|)$.
The digamma integral representation gives the same nonzero-frequency result:
the cosine term contributes $2\pi/(e^{|\omega|}-1)$ to the transform of
$-2\operatorname{Re}\psi(1+ix)$, in addition to $2\pi$ from the delta.
The transform is distributional at zero frequency; this calculation does not
fix a zero-frequency subtraction or justify dropping the finite-sea constant
$2\log h$. The large-frequency limit is $2\pi$, whereas the printed symbol
approaches $4\pi^2$. Four high-precision product checks reproduce the discrepancy.
The original script never tests this symbol. Its subsequent symbolic series
uses $T$, so this particular symbol misprint does not propagate into that series.

\paragraph{Momentum Laplace transform: refuted as printed.}
\Lc\ Let $G(u)=\frac43y(u)^{3/2}$. Since $y(u)\sim u$, its leading Laplace
coefficient must be $\Gamma(5/2)\,4/3=\sqrt\pi$. The normalisation compatible
with the implemented expansion is
\begin{equation}
 \mathcal L G(s)=\sqrt\pi\,T(s)/s^2,
 \label{eq:correctG}
\end{equation}
not $T(s)/s^2$. Without this factor the leading result is
$y(u)\sim\pi^{-1/3}u$, incompatible with the asserted $\delta u=O(\log y)$.
The script silently multiplies every inverse-Laplace term by
\texttt{sp.sqrt(sp.pi)} in \texttt{rhs\_u}; this is not in the displayed theorem.
Also, the printed scale $aP_{\rm rel}/(\sqrt{bh}\,h)$ is inconsistent with
$P_{\rm rel}=4b(\kappa\cosh\kappa-\sinh\kappa)/a$ and
$\kappa^2=yh/b$: the scale giving $4y^{3/2}/3$ is
$a\sqrt b\,P_{\rm rel}/h^{3/2}$. These are algebraic normalisation tests,
not alternative physical conventions.

\paragraph{The printed inverse series is also refuted.}
\Lc\ Put $A=1/(2\pi)$, $c=c_{\rm edge}$, $d=d_0$, $e=e_0$.
Reverting $u=y-A\log y+c+d/y+e/y^2$ necessarily gives
\begin{align}
 y-u&=v_0+v_1/u+v_2/u^2+\cdots,\\
 v_0&=A\log u-c,\qquad v_1=Av_0-d,\nonumber\\
 v_2&=Av_1-\tfrac A2v_0^2+dv_0-e.\label{eq:inverse}
\end{align}
In particular the $u^{-2}(\log u)^2$ coefficient is $-1/(16\pi^3)$.
Phase 2i's equation \texttt{y\_of\_u} has no such term. Numerical inverse
Laplace transforms independently contradict that displayed equation.
The script computes an intermediate inverse series but never compares it to
the printed \texttt{y\_of\_u}; its final reversion can therefore pass.

\paragraph{Coefficients: qualified as a theorem; strong conditional confirmation.}
\Lc\ With \eqref{eq:correctG}, independent numerical inversion supports
$c=(1-\log(8\pi))/(2\pi)$,
$d=1/48-1/(4\pi^2)$ and
$e=(4+6\zeta(3)-\pi^2)/(96\pi^3)$.
Fitting the independent inverse transform without setting logarithmic
coefficients to zero gives
\[
 (d_1,e_2,e_1)\simeq(-2.4\,10^{-21},-2.5\,10^{-15},7.2\,10^{-14}),
 \qquad e_0\simeq0.000451097178279.
\]
This fit uses sixteen $u$ values from $1000$ to $1000(1.5)^{15}$ and includes
all logarithmic terms through $u^{-4}$ in the fitting basis. It is a numerical
check of the corrected transform, not an independent proof of exact zeros.
The original formal SymPy reversion supports exact zeros \emph{conditional}
on its transform input. Independent checks at $u=8,32,128,512,2048,8192$
find the residual after $e/y^2$ smaller than $0.0004/y^3$.

\Ld\ Matching a homogeneous edge solution to a fixed bulk density is a
legitimate way to determine its otherwise free amplitude; it is not intrinsically
an arbitrary convention. Here the projection, zero-frequency matching and
normalisation are asserted rather than worked through, and two printed
normalisations contradict that matching. Hence Phase 2i does not supply the
claimed complete derivation from the finite-$h$ equation. The independent
finite-$h$ test below is additional evidence, not a replacement for it.

\section{Theorem 2: convolution, contact terms and asymptotics}
\paragraph{The double Laplace formula is refuted as printed (\scr, \texttt{wh}, \texttt{finite}).}
\Lc\ In the corrected convention set $B=G'$ and $f=G''$ on the open half-line.
Stirling's limit gives $B(0^+)=\sqrt\pi$, and therefore
\[
 \mathcal L B=\sqrt\pi T/s,\qquad
 \mathcal L f=\sqrt\pi(T-1).
\]
Differentiating the printed convolution
$2\int_0^{u_2}f(u_1-w)B(u_2-w)\,dw$ at $u_1>u_2$ retains the
endpoint $2\sqrt\pi f(u_1-u_2)$. Equivalently, use the causal distribution
$\sqrt\pi\delta+f$, whose Laplace transform is $\sqrt\pi T$.
The compatible full double transform is
\begin{equation}
 2\pi\frac{T(s_1)T(s_2)-1}{s_1+s_2},
 \label{eq:double}
\end{equation}
with its contact term specified, rather than Phase 2i's prefactor $2$.
Thus the printed Laplace pair and the classical normalisation of its
convolution cannot both hold. The original \texttt{verify\_2} does not evaluate
an inverse transform, the two-variable transform, or a finite-$h$ dressed phase.

\paragraph{Convolution after correction: qualified, numerically supported.}
\Lc\ We numerically invert \eqref{eq:correctG} and integrate the first
convolution directly, with 24- and 48-node Gauss quadrature. We compare with
$[\Phi(q+hu,q+hv)-\Phi(q+hu,q)]/h$ from the finite-sea solver at
$(u,v)=(4,1),(8,2),(8,4)$. Both arguments are fixed before $h\to0$;
the corresponding $(Y,r)$ are obtained from the independently computed $y(u)$.
The finite-$h$ grids, extrapolation spread and numerical results are reported
in the accompanying tables and JSON. This is a test of the same two-term
normalisation used in Phases 2f/2h, not an antisymmetrised substitute.

\paragraph{Leading integral and $c_1$: confirmed conditional evaluations (\scr, \texttt{integrals\_data}).}
\Lc\ Independent quadrature verifies
\[
 \frac4\pi\int_0^r\log(s/r)\frac{s^2}{1-s^2}\,ds
 =-\frac4\pi(\chi_2(r)-r)
\]
and the displayed $c_1$ integral against the stated elementary expression,
at five ratios between $0.2$ and $0.8$, with errors below $10^{-54}$.
These are useful conditional integral identities. The balanced logarithmic
coboundary cancels $\log Y$ and $c$ at fixed $t>0$ algebraically.
\Ld\ It does not alone prove a uniform integrated expansion: $Yt^2$ is not
large near $t=0$. The asserted $O(Y^{-2})$ remainder and absence of other
endpoint terms require an endpoint estimate or matched expansion. Replacing
$v$ by $y(v)$ makes the exact measure regular but does not make a large-$y$
expansion uniform at $y=0$. Thus the leading limit and two-loop coefficient
are qualified as conclusions about the Bethe equation, although their formal
integrals and numerical evidence survive.

\paragraph{The claimed four-term reciprocity resolution is refuted as printed.}
\Lc\ A four-term edge-subtracted phase vanishes when either argument is zero.
The asserted symmetric product $2B(u_1)B(u_2)$ does not: $B(0^+)=\sqrt\pi$.
This is a direct boundary-condition contradiction. The earlier Toda dictionary
audit explicitly distinguishes two-term and four-term phases by one-body edge
terms; Phase 2i's remark suppresses that distinction. Its script checks only a
series of the proposed product in $1/Y$, not the endpoint identity or its
finite-$h$ limit. \Ld\ No replacement physical quantum frame is selected here.

\section{Saved Toda data and the finite-$h$ error budget}
\Lc\ The two-term Phase 2f definition is
\[
 R_h=\{\Phi(p_1,p_2)-\Phi(p_1,e_2)-M_1(P_{n2}/a-\Pi_2)\}/h
       -2Y(1-r^2)\log\frac{1-r}{1+r}.
\]
Phase 2h defines its $R_0$ as the $h\to0$ limit of the same expression,
equivalently $\Phi_q(u(Y),u(r^2Y))-\Phi_{\rm cl}(Y,r^2Y)$.
The Phase 2i saved-data subtraction retains this definition in the limit:
the finite-$h$ classical Toda phase has the same leading term.
It is a numerical improvement to the extrapolation, not a new quantum frame.
Our independent high-precision subtraction uses the hyperbolic cross ratio
\[
 A_T=\frac{\sinh^2((\kappa_1-\kappa_2)/2)}
           {\sinh^2((\kappa_1+\kappa_2)/2)},
\]
and the classical momentum and topological term, avoiding the cancellation
in the source script's Hirota-polynomial evaluation.

\Lc\ The original script fixes $c_0$ to its proposed value and computes
\begin{equation}
 \widehat c_1=64[R_0(32)-c_0]-16[R_0(16)-c_0].
 \label{eq:fixedfit}
\end{equation}
This cancels a $Y^{-2}$ term \emph{if} the proposed $c_0$ and a power expansion
are correct. It is different from the old two-point slope
$32[R_0(16)-R_0(32)]$, which estimates $c_0$ freely but neglects curvature.
Consequently its agreement is not an independent simultaneous determination
of $c_0$ and $c_1$. The numerical tables compare both estimators and free
quadratic $1/Y$ fits on $Y=8,16,32$ and $4,8,16,32$.
No fit uses the target $c_1$ as input.

\Ld\ Moving the edge to $q(h)$ does not prove analytic dependence on $h$.
Subtraction of known classical curvature is justified; rejection of every
$h\log h$ correction requires an asymptotic argument about the quantum
finite-sea error. Fit spreads below are sensitivity diagnostics, not
statistical confidence intervals or rigorous truncation bounds. The finite-$h$
coefficient estimates do not resolve exact vanishing of logarithmic coefficients.

\paragraph{Fresh finite-$h$ momentum check (c).}
The grids are $h=0.04,0.02,0.01,0.005$, with respectively
$681,1062,1824,3348$ Gauss nodes. A separate $h=0.01$ solve uses 2189 nodes.
The table compares the four-point cubic intercept to the finest-three-point
quadratic intercept; mesh differences are measured at $h=0.01$.
\begin{center}\small
\begin{tabular}{rrrrr}\toprule
$y$ & Cubic $\delta u$ & Difference from asymptotic & Power-fit spread & Mesh difference\\\midrule
2 & -0.4663923834 & 4.90e-05 & 2.52e-09 & 1.01e-11\\
4 & -0.5757147346 & 4.79e-06 & 1.27e-08 & 1.39e-11\\
8 & -0.6854957971 & 5.52e-07 & 6.62e-08 & 1.68e-11\\
16 & -0.7955382763 & 9.91e-08 & 3.71e-07 & 1.89e-11\\
32 & -0.9057166221 & 3.45e-07 & 2.19e-06 & 2.04e-11\\
64 & -1.0159613570 & 3.48e-06 & 1.29e-05 & 2.15e-11\\
\bottomrule\end{tabular}\end{center}
\Lc\ At $y=16$, extracting $e_0$ from this intercept gives $0.00047646$,
while the power-fit spread propagated by $y^2$ is about $0.000095$.
At $y=32,64$ that amplification is still worse. Thus finite-$h$ data support
$c,d$ but do not independently resolve $e_0$ or exact logarithmic zeros.
The corrected inverse transform provides the sharper conditional test.
\paragraph{Fresh convolution comparison (c).}
\begin{center}\small
\begin{tabular}{rrrrrr}\toprule
$(u,v)$ & $(Y,r)$ & Convolution & BA cubic & Discrepancy & Power-fit spread\\\midrule
$(4,1)$ & $(4.598,0.554)$ & 2.19960638 & 2.19960612 & 2.6e-07 & 5.3e-06\\
$(8,2)$ & $(8.699,0.536)$ & 3.78568886 & 3.78568487 & 4.0e-06 & 4.6e-05\\
$(8,4)$ & $(8.699,0.727)$ & 10.08223894 & 10.08222807 & 1.1e-05 & 1.3e-04\\
\bottomrule\end{tabular}\end{center}
\Lc\ Quadrature changes between 24 and 48 nodes are below $2.6\,10^{-16}$;
these are float-node quadrature comparisons, not 35-digit accuracy claims.
The BA discrepancies are smaller than the power-extrapolation spreads.
The independent mesh comparison changes the three finite-$h$ phases by at most
$3.7\,10^{-11}$ at $h=0.01$.
\paragraph{Saved-data sensitivity after classical subtraction (c).}
All following numbers use the same independently subtracted saved data.
The fixed-$c_0$ estimator is Eq.~\eqref{eq:fixedfit}. ``Free two'' is the
unconstrained intercept/slope estimate at $Y=16,32$.
\begin{center}\small
\begin{tabular}{rlrrr}\toprule
$r$ & $h$ extrapolation & Fixed-$c_0$ $c_1$ & Free two $c_1$ & Target\\\midrule
0.4 & cubic4 & -0.0078316 & -0.0072586 & -0.0077366\\
0.4 & quad3 & -0.0083604 & -0.0070237 & -0.0077366\\
0.4 & quad4 & -0.0102111 & -0.0062013 & -0.0077366\\
0.4 & hlog3 & +0.0065531 & -0.0128715 & -0.0077366\\
0.4 & hlog4 & +0.0233713 & -0.0191761 & -0.0077366\\
0.5 & cubic4 & -0.0215719 & -0.0213581 & -0.0214446\\
0.5 & quad3 & -0.0225280 & -0.0209332 & -0.0214446\\
0.5 & quad4 & -0.0258741 & -0.0194462 & -0.0214446\\
0.5 & hlog3 & +0.0051142 & -0.0317775 & -0.0214446\\
0.5 & hlog4 & +0.0365390 & -0.0435828 & -0.0214446\\
0.6 & cubic4 & -0.0486099 & -0.0488987 & -0.0484360\\
0.6 & quad3 & -0.0501045 & -0.0482344 & -0.0484360\\
0.6 & quad4 & -0.0553354 & -0.0459094 & -0.0484360\\
0.6 & hlog3 & -0.0055121 & -0.0657401 & -0.0484360\\
0.6 & hlog4 & +0.0456838 & -0.0850237 & -0.0484360\\
\bottomrule\end{tabular}\end{center}
\Lc\ ``cubic4'' uses $0.02,0.01,0.005,0.0025$; ``quad3'' and ``hlog3''
use the finest three; ``quad4'' and ``hlog4'' use all four. The logarithmic
basis is $(1,h,h\log h)$. Pure-polynomial alternatives alone shift the
fixed-$c_0$ estimates by more than $2\,10^{-4}$; logarithmic alternatives
still reverse signs. Hence the claim that subtraction resolves the earlier
fit ambiguity is refuted, although the selected cubic comparison passes.
\begin{center}\small
\begin{tabular}{rrrr}\toprule
$r$ & Cubic $h$ fit, free quadratic $1/Y$ (3 points) & Same (4 points) & Fixed-$c_0$ spread\\\midrule
0.4 & -0.0079815 & -0.0078174 & 0.03358\\
0.5 & -0.0213367 & -0.0215472 & 0.06241\\
0.6 & -0.0480635 & -0.0477384 & 0.10102\\
\bottomrule\end{tabular}\end{center}
The last column includes all six saved-data variants, also the coarser
four-point cubic grid $0.04,0.02,0.01,0.005$; it is a sensitivity range.
No statistical or rigorous error interpretation is assigned. \Lc


\section{Theorem 3: correct bookkeeping, false necessity, no physical selection}
\paragraph{The ``iff'' criterion is refuted (\scr, \texttt{kp}).}
\Lc\ Inversion antisymmetry is sufficient for antisymmetry of a ratio-kernel
bilinear under exchange; oddness is unnecessary for that argument.
For example $K(z)=z^2-z^{-2}$ is even and inversion-antisymmetric and gives
antisymmetric functionals. More decisively, all soliton and defect weight
vectors have zero total weight. Adding a constant to $\widetilde K_L$
changes neither $\mathcal F$ nor $\mathcal I_{SS}$, but destroys both
claimed kernel identities. This refutes necessity even while retaining
\emph{both} functional identities universally. A criterion modulo the null
space, with additional domain assumptions, would be a different statement.

\paragraph{Orientation identities are confirmed; physical interpretation qualified.}
\Lc\ We reconstruct the elementary arctangent phonon phase independently,
then integrate $\theta_1\theta'_2/(4\pi)$ on the two open half-lines.
For $(\alpha_1,\alpha_2)=(0.3,0.7)$ all four orientations satisfy
\[
 I^{\rm raw}_{12}=\mathcal I_{SS}[\widetilde K_L]-\epsilon_1\pi/2,
 \qquad J_{12}=\mathcal I_{SS}[\widetilde K_L]
                       -(\epsilon_1-\epsilon_2)\pi/4.
\]
We separately integrate the rational derivative of the defect transmission
phase for $q_D=0.4,0.1,-0.45$, both soliton directions and both defect
orientations. These cover $s<-1$, $|s|<1$, $s>1$, and confirm
\[
 I_{SD}^{\rm raw}=\mathcal I_{SD}[\widetilde K_L]
                         +\epsilon\Delta\theta_D/4.
\]
The maximum discrepancy is below $10^{-58}$ (recorded fully in JSON), with tolerance $10^{-50}$.
Thus the offsets found in the old audit are explained exactly. They have
been moved into explicit orientation constants, not removed from the raw
integrals. \Ld\ This is an algebraic redefinition; it does not calculate
an absolute determinant phase, select $I$ versus $J$, or normalize defect
states. The real-part kernel used on complex arguments is not itself
holomorphic. Holomorphic continuation uses the corresponding branchwise
analytic expression; that distinction must be maintained.

\section{Bootstrap synthesis and defect spectrum}
\paragraph{Crossing confirmed; printed Darboux identity refuted (\scr, \texttt{kp}).}
\Lc\ Independent root sums and the canonical/mode charge shifts check both angle
substitutions. But the printed Darboux defect delay has an extra $\sqrt3$.
With the repository roots $q_{2,a}-q_{2,b}=k_2$ and the independently
reconstructed raw delay
$\Delta x_D=-\epsilon_1\log|(p_{1,b}-q)/(p_{1,a}-q)|/k_1$,
the right side of the printed coboundary exceeds its left side by
\begin{equation}
 2k_2\gamma_1(\sqrt3-1),\qquad
 \gamma_1=-\frac{3\cos2\alpha_1}{\sin\alpha_1\sin3\alpha_1}.
\end{equation}
Replacing its shift $-2\sqrt3\epsilon_1\eta\gamma_1q$ by
$-2\epsilon_1\eta\gamma_1q$ gives the claimed identity in all four tested
orientations. This is an algebraic correction within the existing root
normalisation, not a new physical convention. The original checker never
constructs a defect delay. \Ld\ Crossing and this corrected identity do not
determine a physical fusion constant or supply stable intermediate states.

\paragraph{Breather continuation: qualified.}
\Lc\ Direct continuation of the same three soliton roots used on the real
axis gives $\operatorname{Re}\mathcal I_{SS}=0$ and
$\operatorname{Im}(\Delta P_{\rm can,1}+\Delta P_{\rm can,2})=0$ at both
quoted conjugate pairs, to $10^{-50}$. Branch-point cancellation at $z=1$
is independently checked. The original script instead uses angles
$\alpha\pm\pi/3,\alpha$ in this subsection, different from its real-soliton
$\alpha\pm2\pi/3,\alpha$ roots. The vanishing real part can survive this
change because conjugation and exchange symmetry enforce it; that test alone
does not test the physical root continuation. \Ld\ Pure imaginary phase implies
a positive multiplier only where that phase is finite. The sampled identities
are not a global proof of pole freedom or a quantum quantisation rule.
The shifted Bogdan--Charkina condition is asserted in the text and is absent
from the checker; neither a corrected breather action nor its charge/frame
identification is derived. It remains (d).

\paragraph{Defect correction: confirmed, with a threshold qualification (\scr, \texttt{determinant}).}
\Lc\ Independent differentiation of the nonlinear sewing conditions produces
the linear matrix. Six off-shell determinants confirm the $-2/3$ prefactor.
Filtering the four polynomial roots by $S>1$ confirms the absolute-value
formula for the second growth rate on either side of $s=\sqrt3$.
At the switching point that rate is zero: there are not two strictly positive
growth rates there. The corrected determinant and branch switch were already
established in \texttt{audit\_frames\_defects}; they are successful adoption
of prior (c) results, not new Phase 2i discoveries.

\section{What happened to each finding of the six audits?}
\Lc\ The following is a textual coverage audit, not a revalidation of every
formula in those six documents. ``Relabelled only'' means an outstanding task
is acknowledged but not solved. Grouped entries combine findings with the
same disposition; the explicit omissions are listed separately.
\small
\begin{longtable}{@{}L{9.7cm}L{5.4cm}@{}}
\toprule Prior finding & Phase 2i disposition\\\midrule\endhead
\multicolumn{2}{l}{\textit{Independent Phase 2h audit}}\\
Orientation-dependent $\pi$ discrepancies & Resolved algebraically by explicit offsets.\\
Raw exchange sum and nonzero diagonal & Partially resolved: new kernel antisymmetrises a redefined phase.\\
Oddness alone does not imply antisymmetry & Partially resolved: inversion suffices; new ``iff'' is false.\\
Saved $c_1$ signs change with fit basis & Partially resolved: subtraction helps; robustness is not proved.\\
Leading limit previously built into the loop target & Partially resolved: conditional integral is new; WH derivation incomplete.\\
Classical $(\log Y)^2/Y$ Taylor term & Partially resolved: fixed-$t$ cancellation; endpoint control absent.\\
Separated head-on asymptotics & Resolved in text; original Phase 2i script does not test them.\\
Wrong soliton profile & Resolved in text; original Phase 2i script does not test it.\\
Legendre identity versus physical charge & Relabelled only: field identification remains conditional.\\
\multicolumn{2}{l}{\textit{Charges and action audit}}\\
Regulated co-moving action needs cutoff endpoint & Not addressed by a new regulated calculation.\\
Canonical derivative requires on-shell profile and boundary variation & Relabelled only under general field-theory caveats.\\
Weighted cosine sum is not zero & Not addressed.\\
Printed sign of intermediate $dL_1/d\alpha$ formula & Not addressed.\\
Wrong profile & Resolved by the corrected profile.\\
Static defect has a localized quadratic momentum & Not addressed explicitly; no operator calculation.\\
Fifth-degree boundary potential; bulk tadpole insufficient & Partially resolved: potential acknowledged, counterterm open.\\
Conditional absorption-locus shift, not a charge switch & Not addressed; no quantum sewing/absorption calculation.\\
\multicolumn{2}{l}{\textit{Determinant foundations audit}}\\
Fourth-order domain/completeness differs from cited scalar argument & Relabelled only: full determinant remains open.\\
Wrong co-moving first-derivative sign & Not addressed.\\
Positive potential-field density check & Not addressed by a new check.\\
Orientation-dependent endpoints and exchange constants & Resolved algebraically, physical normalization unresolved.\\
Defect logarithm branch and normalized fusion states & Relabelled only.\\
Amplitude modes have tails; mixed symplectic term does not vanish & Partially resolved: tail term acknowledged; printed Darboux coefficient is wrong.\\
Endpoint Hessian is not field spectral flow or Maslov phase & Not addressed specifically; no spectral-flow calculation.\\
Unit Darboux Jacobian does not fix full functional measure & Relabelled only.\\
Equal localized endpoint frequencies do not prove cancellation & Not addressed; no localized-mode propagation.\\
Stable scattering domain and threshold completeness & Relabelled only; no spectral theorem.\\
\multicolumn{2}{l}{\textit{Frames and defects audit}}\\
Sewing determinant missing factor sixteen & Resolved.\\
Growing branch switch and $s=2$ coalescence & Resolved, except marginal point wording.\\
KP triangle traverses unstable vacua & Relabelled only: physical fusion still open.\\
Edge reference cannot change a ratio slope & Partially resolved: acknowledged; conditional convolution supplies alternative origin.\\
Classical frame does not establish quantum frame & Not resolved; a proposed endpoint product is not the missing dictionary.\\
\multicolumn{2}{l}{\textit{Toda dictionary audit}}\\
Exact edge momentum and dressed-coordinate Jacobian & Not independently derived in Phase 2i.\\
Numerical $q(1+10^{-10})$ offset is nonuniform in $h$ & Not addressed by its solver call.\\
Finite-sea reciprocity has endpoint terms & Partially resolved: invoked, but proposed four-term edge product fails boundary test.\\
Two-term versus four-term normalization matters & Not resolved; the remark conflates their symmetric remainders.\\
Classical frame leaves nonseparable symmetric remainder & Partially resolved: proposes an asymptotic expression without verifying identity.\\
Symmetric part does not fix antisymmetric physical frame & Not addressed by a Ward identity; still open.\\
\multicolumn{2}{l}{\textit{Projected measure audit}}\\
Tail-independent Wronskian determinant and collective volume & Not addressed explicitly; already proved in that audit.\\
Exact collision-time quadrature and conserved boundary flux & Not addressed.\\
Basis-independent projector & Not addressed.\\
Projector violates fixed potential endpoints; fixed-$M$ rank loss & Not addressed beyond generic ``measure open''.\\
Solution tangents need not be temporal determinant zero modes & Not addressed.\\
Volume/endpoint map does not determine propagation phase & Not addressed.\\
Signed momentum derivatives need orientation factor & Not addressed explicitly.\\
Boundary-compatible projected determinant and absolute phase & Relabelled only.\\\bottomrule
\end{longtable}
\normalsize

\section{Script audit, summary and failures recorded}
\Lc\ Running the original script from \texttt{phase2/scripts/} prints
\texttt{ALL CHECKS PASSED} in 101.88 seconds. The symbol, printed inverse
series, transform normalisation, double transform, convolution versus Bethe
data, Darboux delay coboundary, shifted breather condition, and all-six-audit
coverage are not tested as claimed. Its WH check inserts the missing
$\sqrt\pi$; its convolution check evaluates formal coefficient integrals
without constructing that convolution. Its endpoint check expands an assumed
product. Its breather roots differ from the real-root implementation.
The numerical momentum table cites four $h$ values and $y=32$, while the
script actually uses three $h$ values and stops at $y=16$.

\Lc\ The saved-data test is not literally a formula compared with itself:
it uses independent saved solver output. However it fixes $c_0$ to the target
and adopts one extrapolation basis. Its $2\,10^{-4}$ discrepancy is a central-fit
comparison, not a demonstrated total error bar. No evidence was found of
historical tolerance loosening in the available single script snapshot;
the present script cannot establish its editing history. The original
failures list's claim that all prior failures were reproduced is false by
inspection of the tests and the coverage table.

\begin{longtable}{@{}L{6.5cm}L{2.3cm}L{6.3cm}@{}}
\toprule Claim & Verdict & Evidence and scope\\\midrule\endhead
Printed WH symbol, momentum transform and inverse series & Refuted (c) & Product, leading coefficient, reversion; \scr\ \texttt{wh}.\\
$c,d,e$ and absent logarithmic coefficients & Qualified & Strong corrected-transform evidence (c); derivation and exact edge errors incomplete (d).\\
Printed double transform & Refuted (c) & Missing $\pi$ after restoring momentum normalisation.\\
Correctly normalised convolution & Qualified & Independent inverse transform and finite-$h$ comparison (c).\\
Leading residual and elementary $c_1$ integrals & Confirmed (c) & High-precision quadrature of formal integrals.\\
Exact Bethe asymptotics and strict remainder & Qualified & Missing uniform endpoint control (d).\\
Saved-data agreement below $2\,10^{-4}$ & Qualified & Reproduced for one constrained estimator (c); fit sensitivity exceeds that precision.\\
Universal kernel ``iff'' & Refuted (c) & Constant null kernel and even inversion-odd counterexample.\\
SS/SD orientation formulas & Confirmed (c) & Independent raw quadrature, all requested orientations/regimes.\\
Kernel replacement derives physical phase & Qualified & Redefinition with explicit constants; physical normalization open (d).\\
Crossing / printed Darboux coboundary & Confirmed / refuted (c) & Crossing survives; delay shift has extra $\sqrt3$.\\
Breather reality and shifted quantisation & Qualified & Reality checked locally (c); physical shift asserted (d).\\
Defect determinant and branch correction & Confirmed (c) & Independent nonlinear sewing differentiation; marginal-point caveat.\\
Four-term symmetric endpoint product & Refuted (c) & Violates exact edge-subtraction boundary condition.\\
All six audits resolved & Refuted (c) & Finding-by-finding coverage table.\\\bottomrule
\end{longtable}

\paragraph{Failures recorded.}
\begin{enumerate}
\item Prior-claim failures: wrong Fourier symbol, missing momentum-transform
normalisation and physical scale, incompatible printed inverse series, double
transform normalisation, false necessity criterion, extra $\sqrt3$ in the Darboux delay, and four-term endpoint
product; unsupported blanket audit-resolution statement. Checks assert the
existence of these discrepancies rather than asserting the false equalities.
\item The first independent solver run exceeded its inherited fixed momentum
bracket at the largest coarse-$h$ sample. The audit now doubles the bracket
until it encloses the same root; no tolerance or equation was changed.
\item The first direct Darboux assertion failed. Its initial independent
transcription also inverted the raw transmission ratio; this was corrected
against the stated raw delay. The remaining discrepancy is exactly the
printed extra $\sqrt3$, which the final check explicitly asserts. Neither
the target equality nor a tolerance was weakened to hide the failure.
\item One report-assembly command used a repository-relative path while
already inside \texttt{phase2/}, causing \texttt{FileNotFoundError}; rerunning
from the repository root fixed the path without changing any result.
\end{enumerate}

\paragraph{Measured execution and deliverables.}
\Lc\ The final successful resumed checker ran in 106.01 seconds, reusing the
completed finite-$h$ and saved-data sections from the first run. Fresh dense
Bethe solves took 352.46 seconds in total; the original Phase 2i checker took
101.88 seconds. The finite-$h$ cache contains all five runs, with raw phase,
momentum, node and normalization data. Independent convolution quadrature
and data regressions completed before the initially failed Darboux assertion;
the resumed run rechecks the WH diagnostics and all corrected phonon, crossing,
breather, Darboux and sewing assertions. Its final line is
\texttt{ALL CHECKS PASSED}. A normal run of \scr\ reruns every section,
reusing only the dense-solver cache unless \texttt{--refresh} is specified.
No existing implementation was changed; the full repository suite was not
rerun. The new PDF was built twice with \texttt{pdflatex} and all pages
were rendered and inspected. The TeX file includes its numerical tables
and needs no separate table input.

\paragraph{References and bibliography leads.}
\Lc\ No new external reference is cited or claimed read. This audit uses
local project texts and independently checked mathematics. The bibliographic
agent's requested leads remain: the Bogdan--Charkina title/equation designation
differs between Phase 2i and the literature map; the Corrigan--Zambon article
number differs between Phase 2i and the master report. These are unverified
bibliographic leads here, not adjudicated reference corrections. No inference
is drawn from the absence of the git-ignored literature folder.

\paragraph{What is new, and what remains.}
\Lc\ Phase 2i contributes useful conditional higher-order WH reversion
coefficients, an elementary formal $c_1(r)$ integral, improved classical
lattice subtraction, and explicit all-orientation offset identities.
It adopts the earlier determinant and spectrum corrections. \Ld\ Its most
important unresolved problem is a controlled, correctly normalised derivation
linking the finite-sea Bethe phase to a physically defined quantum frame,
including endpoint/contact terms and uniform edge asymptotics. Algebraic
kernel choices cannot substitute for that identification or for the projected
field determinant.
\end{document}
